\documentclass[stu,12pt,floatsintext]{apa7}

\usepackage[american]{babel}
\usepackage{booktabs}

\usepackage{csquotes} 
\usepackage[style=apa,sortcites=true,sorting=nyt,backend=biber]{biblatex}
\DeclareLanguageMapping{american}{american-apa} 
\addbibresource{bibliography.bib}

\usepackage[T1]{fontenc} 
\usepackage{mathptmx} 
\usepackage{hyperref}


\title{Hyperparameter Optimization in Artificial Intelligence (First Draft)}
\author{Paul Beggs}
\duedate{December, 15, 2025}
\affiliation{Hendrix College}
\course{CSCI 335: Artificial Intelligence}
\professor{Dr. Gabriel Ferrer} 

\begin{document}
\maketitle

\renewcommand{\labelenumi}{\theenumi.}

In a nutshell, we work with computers to try to solve problems using programs in the field of computer science. However, it can be common to have instances where we are unable to write a traditional program for a particularly hard problem. These problems can have a search space that is too large (e.g., checkers or chess), or must be generalized so that we can apply a solution to cases that the program has not seen before. Traditionally, these problems are best suited for artificial intelligence, and as these algorithms are capable of handling these problems with sophisticated means. Some of these algorithms employ hyperparameters, or values used before training an algorithm \parencite{weerts2020importance}. In this paper, we will show that as the complexity of training an algorithm increases, so does the need for more complex hyperparameter optimization methods.

We first dive into more detail about hyperparameters and how they are used in machine learning so that artificial intelligence algorithms work as intended. As mentioned before, a hyperparameter is a value that is set by a researcher that is used during training. A common example would be tuning the \(k\) in k-nearest neighbors. Depending on this value, we can set the threshold for a data point being grouped together with similar data, but more on this later. Tuning hyperparameters are commonly done by trial-and-error: a researcher would input a good guess for what they think the optimal hyperparameter would be, then use machine learning to train the algorithm, and review the results to see if the change had any effect on the model's predictions \parencite{bergstra2012random}. However, this manual process can clearly be tedious and full of frustration because finding the best value essentially requires a good guess. This problem is exacerbated when the complexity of the algorithm is increased by adding more hyperparameters, or increasing the time it takes for the algorithm to learn, thus adding additional wait time before the researcher can try another random pair of hyperparameters.

In the domain of artificial intelligence, there is a phenomenon known as ``the curse of dimensionality.'' In which, as we add more data that we want to track, the harder it is to meaningfully identity relationships in the data \parencite{geeksforgeeks_2024_curse}. This problem can sometimes be treated by employing dimensionality reduction techniques, but we will not be exploring that in this paper. Now, for the times when we cannot reduce the amount of dimensions that we have for our data, such as training a neural network, we must use more sophisticated means to optimize the hyperparameters. 

Thus, we introduce four hyperparameter optimizers: grid search, random search, Bayesian optimization, and simulated annealing. We will then investigate three distinct problems in artificial intelligence: handwriting analysis, sentiment analysis, and image Classification, and how each optimizer can be applied to the problem. Depending on the application, we will determine if one or more optimizer is best suited for the problem, and walk through why that is the case. 

\section{Hyperparameter Optimizers}

\subsection{Grid Search}

One of the most common and straightforward hyperparameter optimizers is grid search. This optimizer works by doing an exhaustive search over a domain and trying every possible combination of hyperparameters together. The search starts by having a researcher manually define points that are equally spaced across a domain. Then, they run loops for each hyperparameter choice, and compare the accuracy between the predicted values against the real values (often called ground truth) for every combination. The hyperparameters that result in the best accuracy. Note that there are many metrics that measure accuracy, like root mean square error or negative log likelihood, but we will not be getting to that in this paper, and will instead assume a generic accuracy metric.

Since hyperparameters can take the form of any type of variable, they range from having infinite possibilities (numerical), to single digit possibilities (categorical). For example, continuous hyperparameters could be \(c\) or \(\gamma\) (regularization hyperparameters), or linear or RBF (kernel type) for support vector machines \parencite{scikitlearn_2019_sklearnsvmsvc}. On one hand, we would have to search an arbitrarily large number of values for \(c\) and \(\gamma\), but on the other, we only need to search five kernels. Hence, categorical variables are where grid search shows the most promise with exhaustive search. 

From grid search's description, it is not hard to imagine that as the number of hyperparameters increases, so does the search time of the space. This is due to exponential time complexity of running loops for every hyperparameter. That is, if we have \(k\) hyperparameters, and we decide to test \(n\) distinct values, the total number of evaluations is \(n^{k}\). To put that into practice, consider the following example: we have 4 parameters with 10 values, which directly translates to \(10^{4} =\) 10,000 training runs. If each run takes 5 minutes, the search would take over a month. 

While running 10,000 training runs can be a major roadblock for most applications, grid search has an extremely beneficial property of being ``embarrassingly parallel.'' Because the results of each set of hyperparameters are independent of each other (i.e., set one's results do not affect set two), we can run multiple hyperparameter configurations at the same time across multiple systems \parencite{bergstra2012random}. Essentially, if we have 10 GPUs, we can train 10 models at the same time, thus reducing the wait time by 90\%. This makes grid search viable for organizations that have massive computational resources, even when the time complexity is exponential. However, other organizations may prefer an option that is just as good as grid search, but does not suffer from the curse of dimensionality because of its versatility. They could consider another optimizer, random search. 

\subsection{Random Search}

Random search, like grid search, has a researcher choose candidate hyperparameter values across a grid in an effort to get the highest success for a model's prediction. However, it differs in how these points are chosen, and does not perform exhaustive searches over the whole domain. Instead, random search picks hyperparameters based on a uniform distribution \parencite{pilario2020kernel}. Given this random sampling of parameters, if we use the same number of points as grid search does, random search finds the hyperparameters that have the highest success for the model \parencite{bergstra2012random}.

The performance gains stem from the fact that random search does not use a fixed grid structure. We initialize the search by picking candidate points from our model's hyperparameter sets, but then we utilize a uniform distribution to give every possible hyperparameter within a preset range an equally likely probability of being picked \parencite{zabinsky2009random}. If our success metric reports higher values for a particular set of hyperparameters, we add them to a list of promising candidates. Otherwise, we continue this process until we reach a cutoff threshold for the success metric or a budget cap.

To show the advantage of random search over grid search, consider a scenario where we can utilize ``low effective dimensionality'' \parencite{bergstra2012random}: while tuning a neural network, we find that only two hyperparameters are of importance: learning rate and momentum. The learning rate controls the size of each step taken toward a minimum, whereas momentum influences how quickly the gradient accelerates in that direction \parencite{orr_1999_momentum}. In this example, we will assume that the learning rate is the more important parameter over momentum. If we were to use a \(3 \times 3\) grid search, we would do nine tests, but because we reused values that we already tried, we essentially only used three distinct values for the learning rate. We wasted six trials of trying those same learning rate values with momentum values that were not as important. If we instead used random search for the same nine trials, we would try nine distinct values for the learning rate, thereby casting a wider net for both hyperparameters. This example shows that random search is much more efficient in high dimensional spaces as samples hyperparameters more effectively.

We have explored how random search can be better than grid search, but that does not mean random search is not without its limitations. Since it, like grid search, each treat sampled hyperparameters as independent of one another, it does not ``learn'' from its past exploration. That is, if one hyperparameter yields a great accuracy score, it does not focus around that region. Instead, it will randomly sample another hyperparameter. The lack of learning from experiences allows for more sophisticated methods, like Bayesian optimization, which use them as a fundamental part of their algorithm.  

\subsection{Bayesian Optimization}

Unlike grid search and random search, which do not take previous data into account, the Bayesian optimization process learns from previous data and incorporates them into a probabilistic model of an ``objective function''\footnote{We can also think of the objective function as our generic success metric that we have been using.} \parencite{frazier2018tutorial}. The objective function is generally considered to be continuous, but it does not have linearity, concavity, and ``first- or second-order derivatives'' \parencite{frazier2018tutorial}. These qualities make the objective function into a ``black box problem,'' where we essentially do not know what its qualities are, and where every evaluation is computationally expensive \parencite{WU201926}. So, we are trying to predict hyperparameters using a probabilistic model for a function that we know very little about. Thus, with the Bayesian process, we can make educated guesses that perform better than other optimization processes \parencite{WU201926}.

Predicting the objective function is similar to how Bayesian statisticians model posterior distributions. They start with a ``hunch'' of what they think the distribution may look like, called the prior distribution. This hunch comes from previous experiences with the objective function, and does not have to be based on fact \parencite{WU201926}. Instead, a researcher can use their opinion for what they think the prior distribution would look like.  Then, they introduce new data to the model and calculate a new posterior distribution which gives them information of where to sample next. Through this process of educated guessing, we can approximate where the best point of the objective function is, of which corresponds to the optimal set of hyperparameters that maximizes the objective function.

After the posterior distribution is updated, the algorithm uses an acquisition function to determine which hyperparameter configuration to test next \parencite{WU201926}. There are many acquisition functions available, but the one we will discuss here is the upper-confidence bound (UCB) acquisition function as it is mainly concerned with balancing two competing strategies: exploration and exploitation. Exploration is the act of exploring areas of high variance, where the function could be much higher or lower. By sampling these regions of high variance, we are exploring possible points that could be better than what the current point has near it. Exploitation, on the other hand, only focuses on finding high probabilities of success, or the largest point on the graph. With UCB, the balance between exploration and exploitation is controlled using a hyperparameter \(k\). The higher \(k\) is, the more it values exploration; the lower it is, the more it values exploitation.

\subsection{Simulated Annealing}

The simulated annealing algorithm stems directly from the physical process of material annealing. That is, heating a material to a high temperature, then cool the material with a tightly controlled process that allows the atoms to be arranged in an optimal pattern \parencite{rutenbar2002simulated}. For hyperparameter optimization, we start by exploring the space by trying any hyperparameter configuration, regardless of improvements for our success metric. Then, the algorithm becomes more selective about the hyperparameters that it chooses, and it eventually settles with what it found to be the best configuration. This optimization process directly matches the process of physical annealing: start with a high temperature, then slowly work your way down until the atoms reach an optimal orientation. 

One of the first simulated annealing algorithms developed is the ``Metropolis algorithm,'' which dictates how the system explores and eventually settles \parencite{rutenbar2002simulated}. It works by initializing a temperature \(T\) and a stopping criterion, \(M\), for which \(T\) slowly decrements, and a counter enumerates until it reaches \(M\). For each loop, we randomly pick a hyperparameter configuration close to the current position, and check the change in the system's ``energy,'' denoted as \(\Delta E\). This energy value is related to the model's error rate, which should be minimized.

If \(\Delta E\) is negative (the new error is lower), we approve of the change, and move toward a minimum. However, if \(\Delta E\) is positive (the new error is higher), the algorithm does not immediately reject the move. Instead, it accepts this upward move based on the Boltzmann probability \(P = e^{\Delta E/T}\). Then, when \(T\) is high, the probability \(P\) is close to 1, allowing the algorithm to accept almost any move. This chaotic behavior allows the search to jump out of local minima—shallow valleys that would trap a greedy algorithm. As \(T\) decreases, \(P\) drops, and the algorithm becomes increasingly selective. Then, either \(T\) hits a cutoff, or our stopping criteria, \(M\), is met, and it settles into the lowest error state found.


\section{Handwriting Classification}

Handwriting classification with artificial intelligence starts by using data (a drawing of the letter `A') and labeling it (`A'). Then, we use that labeled data to train an algorithm to recognize the drawing of the letter `A' and classify it correctly. We can evaluate how good algorithm is by giving it new drawings and asking it to classify it as a specific letter. In general, we would use a cross-validation set, which partitions all data into training and evaluation sets, where all sets do not have any overlap.

To classify data, we could try a specific algorithm, k-nearest neighbors, which works by using distance function to measure the similarity between a new drawing and our labeled data. Suppose that the drawings are composed on a \(40 \times 40\) drawing grid, where the pixels can be either on or off. This means the model must determine which known examples are most similar to this new grid. This can be done by calculating the distance between the input value and every element in the training data. Then, these distance-label pairs are inserted into a priority queue, which has the property of automatically ordering from nearest to farthest. Thus, allowing us to access the most relevant data points without sorting the entire dataset manually. Once the algorithm organizes the distances, we can take the top \(k\) elements from the priority queue and classify them as the nearest neighbors, and add them to a histogram. Now, we can take a vote to get the final classification by counting the frequency of each label among the \(k\) neighbors. The label with the highest count is the predicted classification. 

As we can see, the accuracy of classifying data relies heavily upon what the value of the hyperparameter \(k\) is. If we set \(k\) too high, the model will misclassify a letter as another, but if we set it too low, then the model will not be able to generalize properly, this is known as over-fitting. So, we're left with a hyperparameter optimization problem.

\subsection{Applying Optimizers}

For this problem, we must optimize hyperparameter \(k\) to maximize the accuracy of the histogram's voting process. Since \(k\) is the only hyperparameter that we are testing, grid search is a viable choice as we can define a range for \(k\), say 1 to 20, and exhaustively check each integer in that range to find the highest accuracy. We would not want to try the other hyperparameter optimizers, as we do not have enough hyperparameters to justify the choice. If we had introduced another hyperparameter like a different distance metric, or threshold values for pixel transparency, then the search space would grow, and something like random search would be viable. 

\section{Sentiment Analysis}

Within natural language processing, we often want to understand the sentiment (how they are feeling) of someone who bought a product, or is talking to a friend. Our goal is to see if they are positively or negatively affected by the product. With that information, we could recommend similar products in the form of ads or suggest a different product to the customer, for example. To analyze sentiment, we use a similar process to handwriting classification: we gather data (e.g., reviews on Amazon), and give it as positive or negative. There are many algorithms that handle sentiment analysis; in fact, we could use the k-nearest neighbors algorithm again, but this process is generally done with a support vector machine, or SVM. 

If we want to classify data with an SVM, we must contend with multiple hyperparameters. The SVM has 2 sets of hyperparameters: kernel and regularization. The four most common kernels for SVM are linear, polynomial, radial basis function (RBF), and sigmoid. Then, the regularization values, \(c\) and \(\gamma\), are applied to the kernels in different ways. The hyperparameter \(c\) is only used for the polynomial and sigmoid kernels, whereas \(\gamma\) is used in every kernel except linear \parencite{jain_2024_svm}. Since we usually work with just two sentiments, it is common to pick the linear kernel over the others, but we cannot know for certain that for every dataset, that this kernel is the best \parencite{reddy_2023_sentiment}. This supplies another problem that we can pursue hyperparameter optimizers with.

\subsection{Applying Optimizers}

When looking to hyperparameter optimizers, we might first look to grid or random search. As previously noted, grid search can perform the best when dealing with categorical variables within a limited domain, such as our four kernel choices. So, if we were only selecting a kernel, grid search would be the best because we could test all four options in relative quickness. However, the inclusion of the regularization parameters \(c\) and \(\gamma\), make grid and random search non-optimal. Instead, we turn to simulated annealing to navigate the hybrid search space. As mentioned earlier, simulated annealing works by choosing a new configuration that is close to the current one. Since we are using continuous parameters, this is relatively straightforward: we just pick a value for \(c\) or \(\gamma\) that is slightly higher or lower than what we currently have. For the kernel type, being close just means we arbitrarily change to one of the three options. This flexibility allows us to explore the discrete choice of kernels, while also fine-tuning the continuous regularization parameters.  

\section{Image Classification}

In the previous section about handwriting classification, we greatly simplified the classification problem by restricting the domain to only letters with only binary values for the pixels. If we want to employ a generic image classifier that is able to handle any image, we must look towards convolutional neural networks (CNN) to handle the increased amount of details. That is, we now have an enormous amount of pictures each containing arrays of pixel values for the red, green, and blue values. 

For this problem, we can leverage existing trained models, like VGG16, which contains 138 million parameters. This CNN was trained with the ImageNet database, which has over 14 million labeled training images \parencite{gabrielferrer_2025_csci}. We can configure this model so that we don't affect its trained parameters. Then, we can build another neural network to the end of the model to perform specific image classification of our dataset (e.g., cats and dogs). However, we see that this addition introduces a ton of hyperparameters. We must find an optimal rescaling size to cut down on the neural network's training time. In addition, we must decide how many layers we are adding, and for those layers, what activation function to use for them. In short, there are many hyperparameters that we need to optimize, and that is not even taking the training into account (e.g., learning rate, loss function, optimizer, etc.). 

\subsection{Applying Optimizers}

This is a prime problem for Bayesian optimization. We are handling a neural network that takes time to train, and (depending on available resources) is energy expensive. We need the optimizer to be efficient and effective at picking the optimal hyperparameters. As discussed earlier, the optimizer treats the model's accuracy as a ``black box'' function that is expensive to evaluate. So, we will build a probabilistic model of how the layers and activation functions affect accuracy. In turn, the algorithm will learn from each run, and incrementally get better. In this situation, every training run matters, so we cannot use exhaustive measures by any means. 

\section{Conclusion}

Through this paper, we have explored 4 distinct hyperparameter optimizers: grid and random search, Bayesian optimization, and simulated annealing. We discussed how grid search is best when we have relatively few hyperparameters, and when we have resources to spare to parallelize the search. Then, we looked towards random search, which had a lot of the same properties as grid search, but performs better in general. This is because random search has low effective dimensionality: since it samples hyperparameters uniformly, it has a wider breadth over important hyperparameters. This is because we showed that grid search wastes 6 tests for a \(3 \times 3\) grid. Then, we moved to Bayesian optimization, a complicated algorithm that uses a probabilistic model for an objective function by employing prior distributions, past data, and acquisition functions. A specific acquisition function that we looked into, upper-confidence bound (UCB) balances exploration versus exploitation, to make our search as effective as possible. We finished off with simulated annealing, an algorithm that is based on the physical process of annealing material. This algorithm excels at avoiding minimums and maintaining a wide breadth of the search space. 

After we introduced each algorithm, we then set up common problems that come up in artificial intelligence, and walked through each problem by considering how each optimizer can be applied. For the first problem, we looked into handwriting classification and found that since we are using k-nearest neighbors, which only has one hyperparameter, that grid search would be best. This is because we could define a range of possible \(k\) values and try each one without having to worry about the curse of dimensionality. Then, we moved onto sentiment analysis with support vector machines. We decided that grid and random search would not be adequate, as the variables are continuous. This left us with simulated annealing, as we showed that properties to explore the global space were the most adequate. This left us with the hardest problem, image classification. We went through an example where we could use VGG16, a powerful convolutional neural network that was pretrained on millions of images. We wanted to build a neural network on top of this one, so we had to contend with at least 6 hyperparameters. This made Bayesian optimization the most appealing, and the most cost-effective option. 

Therefore, we have shown through our exploration of hyperparameter optimizers and specific problems that as the complexity of training an algorithm increases, so does the need for complex hyperparameter optimization methods.

\printbibliography

\end{document}
